% Put your abstract here

\begin{abstractwithpageno}	
International School Manila (ISM) maintains large volumes of academic and
administrative reports whose authenticity and integrity are critical for
accreditation, audits, and external verification. At present, reports are
stored in centralized systems without cryptographic binding between the files
and their metadata, leaving room for undetected modification, incomplete audit
trails, and slow, IT-dependent verification. This project proposes
\textbf{VERACISM} (Verifiable Academic Records for ISM), a hybrid platform that
anchors finalized reports to content-addressed, decentralized storage while
keeping identity, access control, and operational monitoring under ISM
governance.

In VERACISM, reports generated by existing ISM applications are uploaded
through a \texttt{.NET~8} Web API, scanned using ClamAV and YARA, validated by
MIME type and size, hashed with SHA-256, and written to IPFS with
Filecoin-backed deals via Lighthouse. Each accepted file receives a unique
Content Identifier (CID), which is stored together with report metadata and
immutable audit events in MSSQL under strict RBAC/ABAC rules. A web-based
console supports tenant and integrator administration, while a public
verification portal allows internal and external parties to confirm report
integrity using CIDs or QR codes, optionally cross-checking storage evidence in
public explorers.

The project follows a DevSecOps lifecycle, integrating secure design,
automated checks, and scheduled VAPT aligned with OWASP Top~10 risks. Planned
evaluation combines task-based user testing for administrators, tenant offices,
integrators, and verifiers with structured security testing and defined
acceptance criteria. By replacing implicit trust in storage platforms with
cryptographic proof and transparent audit trails, VERACISM aims to strengthen
ISM’s record governance, reduce the overhead of verification workflows, and
provide a reusable pattern for tamper-evident academic report storage.
\end{abstractwithpageno}
